\documentclass[
]{ceurart}

\sloppy

\usepackage{booktabs}
\usepackage{graphicx}
\usepackage{listings}
\usepackage{multirow}
\usepackage{url}
\lstset{breaklines=true}

\begin{document}

\copyrightyear{2026}
\copyrightclause{Copyright for this paper by its authors.
Use permitted under Creative Commons License Attribution 4.0
International (CC BY 4.0).}

\conference{CLEF 2026 Working Notes, September 21--24, 2026, Jena, Germany}

\title{Hybrid Dense-Sparse Retrieval and Re-Ranking for Multilingual Resume Retrieval from Job Descriptions}
\tnotemark[1]
\subtitle{Working Notes for TalentCLEF 2026 Task A}
\tnotetext[1]{Search runs and code are available at \url{https://github.com/alonsopg/talentclef2026-taskA}.}

\author{Alonso Palomino}[email=mail@alonsopg.com]
\cormark[1]
\cortext[1]{Corresponding author. Email: \href{mailto:mail@alonsopg.com}{mail@alonsopg.com}.}

\begin{abstract}
This paper describes a submission to TalentCLEF 2026 Task A, the contextualized
job-person matching task. The objective is to rank candidate resumes for a given
job description in a multilingual setting. The study evaluates a two-stage retrieval architecture
based on JobBERT as the first-stage dense retriever and compares two re-ranking strategies:
a sparse lexical re-ranker based on SPLADE and a late-interaction re-ranker based on
multilingual ColBERT. On the official development split, the strongest overall system
is obtained with JobBERT followed by SPLADE-based reciprocal-rank-fusion re-ranking.
ColBERT remains competitive in Spanish and provides a more interaction-rich alternative
for long-form resume matching. The best dense-sparse system is used as the primary run.
\end{abstract}

\begin{keywords}
TalentCLEF \sep
job-person matching \sep
resume retrieval \sep
dense retrieval \sep
SPLADE \sep
ColBERT \sep
multilingual information retrieval
\end{keywords}

\maketitle

\section{Introduction}

TalentCLEF 2026 Task A focuses on contextualized job-person matching, where each
query is a full job description and the objective is to retrieve and rank the most
relevant candidate resumes. This setting is particularly challenging because job
descriptions and resumes are long, structured, multilingual, and information-dense.
As a result, effective systems need to capture both broad semantic compatibility and
precise lexical evidence such as skills, tools, certifications, and role-specific terms.

Viewed more broadly, retrieval over skill requirements and competence descriptions
extends beyond job-to-candidate matching. Closely related problems also arise in
educational settings, where systems must identify assessment items that measure
targeted vocational skills. Prior work has shown that semantic retrieval and
re-ranking can support vocational exam assembly by retrieving relevant test items
from validated item banks and by improving ranking quality under practical constraints
such as bias control \cite{palomino2024edtec,palomino2025bias}. This broader perspective
highlights that skill-aware retrieval methods are useful wherever textual evidence
must be matched to explicit competency requirements.

For this reason, a two-stage architecture is adopted. In the first stage,
JobBERT is used as a dense retriever to generate a high-quality candidate set.
In the second stage, the candidates are re-ranked using either a sparse lexical
signal from SPLADE or a late-interaction signal from multilingual ColBERT.
This design makes it possible to study two complementary ways of strengthening
the first-stage semantic ranking.

The paper addresses the following research question: can a two-stage multilingual
retrieval pipeline improve resume retrieval from job descriptions over a strong
JobBERT baseline by adding a second-stage re-ranker with either sparse lexical
matching or late interaction? The working hypothesis is that second-stage re-ranking
will improve retrieval effectiveness over JobBERT alone, and that sparse re-ranking
with SPLADE will be more robust overall than zero-shot multilingual ColBERT because
explicit lexical evidence such as skills, tools, and credentials remains central to
resume relevance.

The main finding is that re-ranking helps, but the choice of second-stage model matters:
SPLADE is the most robust overall, while multilingual ColBERT remains competitive,
especially in Spanish.

\section{Task Description}

TalentCLEF 2026 Task A evaluates contextualized job-person matching in a multilingual
retrieval setting. Given a job description, the system must rank candidate resumes
according to their relevance to the job requirements. The submission focuses on the monolingual
\texttt{en-en} and \texttt{es-es} tracks.

The official submission format follows the TREC run-file convention, with one run file
per language pair. Submissions are evaluated with standard information retrieval
metrics, with Mean Average Precision (MAP) as the main ranking criterion and nDCG
reported as an additional metric \cite{talentclef2026eval,talentclef2026format}.

\section{Data and Preprocessing}

\subsection{Official Data}

This work uses the official TalentCLEF 2026 Task A development and test sets
\cite{talentclef2026task}. The development data contains 10 queries and 472 resumes
for each of English and Spanish. The test data contains 40 queries and 476 resumes
for each of the two languages.

\subsection{Resume Text Construction}

The final systems do not use the earlier ESCO-based or knowledge-graph-based re-ranking
components explored in preliminary experiments. Instead, a text representation is built
for each resume by extracting and concatenating the main resume sections, including
profile, experience, skills, education, certifications, and languages. Both resumes
and job descriptions are normalized before retrieval.

This preprocessing design keeps the retrieval pipeline simple and reproducible while
preserving the information most likely to matter for candidate ranking.

\section{Method}

\subsection{First-Stage Dense Retrieval with JobBERT}

As the first-stage retriever, this submission uses \texttt{TechWolf/JobBERT-v3}, a dense encoder
from the JobBERT family developed for labor-market matching tasks \cite{jobbert}.
JobBERT encodes both job descriptions and resume representations into a shared
embedding space. Candidate resumes are then ranked by similarity to the job-description
embedding.

Conceptually, JobBERT is used here as a domain-adapted bi-encoder. The job description
and each resume are encoded independently into dense vectors, which makes retrieval
efficient because resume embeddings can be computed in advance and compared to the
query embedding with a vector-similarity operation. Its main strength is broad semantic
matching when compatible roles or skills are expressed with different wording.

\subsection{Sparse Candidate Re-Ranking with SPLADE}

To complement dense retrieval with stronger lexical matching, JobBERT candidates are re-ranked
using SPLADE \cite{splade,spladev2}. SPLADE produces sparse lexical
representations that retain neural expressiveness while remaining sensitive to exact
or near-exact term overlap.

Unlike a dense bi-encoder, SPLADE maps text into a high-dimensional sparse vector
over the vocabulary. Important terms receive larger weights, and the model can also
activate related expansion terms learned by the transformer. This makes SPLADE useful
for re-ranking in settings where exact evidence such as software names, certifications,
or specific skills can strongly influence relevance.

SPLADE is applied over the JobBERT candidate pool and the two ranking signals are combined
through reciprocal-rank fusion over the candidate set. This yields the strongest
overall system on development.

\subsection{Late-Interaction Re-Ranking with ColBERT}

As a more interaction-rich alternative, the study evaluates a pretrained multilingual
ColBERT re-ranker using \texttt{jinaai/jina-colbert-v2} \cite{colbert,jinacolbert}.
Unlike bi-encoder retrieval, ColBERT preserves token-level interactions between
query and document representations and is therefore appealing for long-form resume
ranking, where local skill- and role-level matches may be important.

More specifically, ColBERT encodes query tokens and document tokens separately and
delays most of the interaction until scoring time. For each query token, the model
finds the best matching document token and then aggregates those token-level matches
into a final relevance score. This late-interaction design is more expensive than
single-vector retrieval, but can capture fine-grained alignments important in resume ranking.

ColBERT is applied as a second-stage re-ranker over the top JobBERT candidates. This
variant is not the strongest overall system, but it is highly competitive in Spanish
and provides a stronger innovation story than a purely dense-sparse hybrid.

\section{Experimental Setup}

\subsection{Systems Compared}

The following systems are compared on the official development split:

\begin{itemize}
  \item \textbf{BM25}: lexical baseline.
  \item \textbf{SPLADE}: sparse neural baseline.
  \item \textbf{JobBERT}: dense first-stage retrieval baseline.
  \item \textbf{JobBERT\texorpdfstring{$\rightarrow$}{->}SPLADE}: candidate re-ranking with sparse lexical evidence.
  \item \textbf{JobBERT\texorpdfstring{$\rightarrow$}{->}ColBERT}: candidate re-ranking with late interaction.
\end{itemize}

Augmentation and compact-query variants were also explored in preliminary experiments,
but these did not outperform the main re-ranking systems and are therefore treated as
secondary ablations.

\subsection{Implementation}

JobBERT retrieval is implemented with SentenceTransformers \cite{reimers2019sentencebert}.
PyTerrier is used for sparse retrieval experiments, and the SPLADE re-ranker is
implemented through the \texttt{pyt\_splade} stack. The late-interaction re-ranker is
implemented with a pretrained multilingual ColBERT model through the PyLate stack.

\section{Development Results}

Table~\ref{tab:dev-results} summarizes the main development results used for model
selection.

\begin{table}[t]
  \centering
  \caption{Main development results on TalentCLEF 2026 Task A.}
  \label{tab:dev-results}
  \begin{tabular}{llcc}
    \toprule
    Language & System & MAP@50 & nDCG@50 \\
    \midrule
    en & JobBERT & 0.7744 & 0.9245 \\
    en & JobBERT$\rightarrow$SPLADE\_RRF & 0.7808 & 0.9294 \\
    en & JobBERT$\rightarrow$ColBERT & 0.7728 & 0.9225 \\
    \midrule
    es & JobBERT & 0.7299 & 0.8962 \\
    es & JobBERT$\rightarrow$SPLADE\_RRF & 0.7350 & 0.9034 \\
    es & JobBERT$\rightarrow$ColBERT & 0.7348 & 0.9086 \\
    \bottomrule
  \end{tabular}
\end{table}

The results show that JobBERT is already a strong dense baseline. Re-ranking with
SPLADE yields the best overall MAP in both English and Spanish. The ColBERT re-ranker
does not surpass the SPLADE re-ranker overall, but in Spanish it is nearly tied in MAP
and slightly stronger in nDCG. This makes ColBERT an attractive additional variant for
analysis and submission.

\section{Analysis}

The experiments suggest that the strongest gains come from combining dense semantic
retrieval with a complementary re-ranking signal. SPLADE improves on the dense retriever
by reinforcing exact lexical evidence, which remains highly relevant in resume matching.
This is especially useful when important job requirements are expressed through specific
skills, technologies, tools, or credentials.

ColBERT offers a different advantage. By preserving token-level interactions, it is
better aligned with the long-document nature of job descriptions and resumes than a
plain single-vector model. Although it was not the best system overall on development,
its strong Spanish performance suggests that late interaction is a viable direction for
future work in multilingual job-person matching.

Several additional variants did not help. Standalone sparse retrieval was clearly
weaker than JobBERT, long query augmentations were often neutral because they were
truncated by the encoder, and compact rewrite variants were not robust, especially in
Spanish. These negative results suggest that extra complexity is useful only when the
added signal is both informative and compatible with the retrieval model's input limits.

\section{Limitations}

This work has several limitations. The public development set is small, which makes
model selection fragile; the public data does not expose gender labels, which limits
direct fairness analysis; and ColBERT was not fine-tuned because public supervision
is limited and zero-shot re-ranking was already computationally demanding.

\section{Conclusion}

A two-stage retrieval framework for TalentCLEF 2026 Task A was presented based on
JobBERT first-stage retrieval and alternative re-ranking strategies. The best
development results were obtained with a hybrid dense-sparse re-ranking approach using
SPLADE, while multilingual ColBERT provided a competitive and more distinctive
late-interaction alternative, especially in Spanish. These findings highlight the
value of combining domain-adapted dense retrieval with complementary re-ranking signals
for contextualized multilingual job-person matching. In relation to the research
question, the results indicate that second-stage re-ranking can improve over a strong
JobBERT baseline in multilingual resume retrieval from job descriptions. The working
hypothesis is therefore only partially confirmed: re-ranking was beneficial overall,
and SPLADE was the most robust approach, while zero-shot multilingual ColBERT remained
competitive but did not consistently surpass the dense-sparse hybrid. Overall, strong
semantic retrieval remains the foundation of this task, and its best results come from
adding a complementary lexical or late-interaction signal rather than from increasing
complexity indiscriminately.

\section*{Declaration on Generative AI}
During the preparation of this work, the author used OpenAI GPT-5.4/Codex in
order to help structure the manuscript, convert a paper outline into a LaTeX draft,
and revise wording. After using this tool, the author reviewed and edited the
content as needed and takes full responsibility for the publication's content.

\bibliography{references}

\end{document}
